\chapter{Proportions et évolutions}
\section*{Exercice 17.1 -- Une proportion}
$N=54/0{,}18=300$. Le calcul inverse utilise la proportion comme coefficient, et non comme effectif.

\section*{Exercice 17.2 -- Proportion de proportion}
$0{,}40\times0{,}30=0{,}12$, soit $12\,\%$. Le deuxième pourcentage porte sur le sous-groupe des usagers du bus.

\section*{Exercice 17.3 -- Successions}
Le coefficient global est $1{,}10\times0{,}85=0{,}935$. Le taux global est $-0{,}065$, soit une baisse de $6{,}5\,\%$.

\section*{Exercice 17.4 -- Réciproque}
Le coefficient initial est $0{,}7$. Il faut multiplier ensuite par $1/0{,}7=10/7$, soit une hausse de $3/7\simeq42{,}86\,\%$. Une hausse de $30\,\%$ ne suffirait pas.

\section*{Exercice 17.5 -- Points ou pourcentage}
Elle augmente de $3$ points de pourcentage. La hausse relative est $(15-12)/12=1/4=25\,\%$. Les deux réponses ne désignent pas la même grandeur.

\section*{Exercice 17.6 -- Pondérer}
Les effectifs oui sont $10$ et $30$, soit $40$ sur $100$ : $40\,\%$. La moyenne simple des deux pourcentages donnerait $37{,}5\,\%$, car elle attribuerait à tort le même poids aux deux groupes.
